\chapter{Спектральная проверка жёсткости майорановского спаривания}

Предыдущий раздел установил геометрический порог
\[
 \lambda v^2>\frac{\pi^2}{L^2}.
\]
Теперь требуется выяснить, следует ли левая часть из уже принятого
спектрального действия.

\section{Разложение спектральной функции}

Пусть \(a\) --- амплитуда майорановской деформации нулевой моды. В
одномерном контрольном секторе
\begin{equation}
 f(a^2)
 =f(0)+f'(0)a^2+\frac12f''(0)a^4+\cdots .
\end{equation}
Сравнение с потенциалом Гинзбурга--Ландау даёт
\begin{equation}
 c_2=f'(0),\qquad
 c_4=\frac12f''(0),\qquad
 \lambda=f''(0),\qquad
 v^2=-\frac{f'(0)}{f''(0)}.
\end{equation}
Поэтому пороговая комбинация имеет особенно простой вид
\begin{equation}
 \boxed{\lambda v^2=-f'(0)}.
\end{equation}
Это равенство относится к приведённому одномодовому сектору с канонической
нормировкой производного члена. Полное действие обязано независимо вывести
эту нормировку.

\section{Проверка принятого набора функций}

Для уже использованных в программе спектральных функций получаем:
\[
\begin{array}{c|c|c}
 f(x) & \lambda v^2 & \text{состояние при }L=\pi\\
\hline
 x & \text{устойчивый минимум не возникает} & \text{нормальная фаза}\\
 e^{-x} & 1 & \text{критическая точка}\\
 e^{-2x} & 2 & \text{конденсат}\\
 (1+x)^{-3} & 3 & \text{конденсат}\\
 (1+x)^{-4} & 4 & \text{конденсат}
\end{array}
\]
Следовательно, различные допустимые спектральные функции дают различные
фазы. Сам носитель не выбирает между ними.

Каноническая норма пространства конфигураций
\[
 23+\pi^{-1}
\]
задаёт положительную кинематическую метрику коллективной деформации, но не
определяет знак \(f'(0)\). Поэтому она не может сама по себе вызвать
неустойчивость нормальной фазы.

\begin{nogocriterion}[Неопределённость спектральной функции]
Пока функция \(f\), её масштаб и относительная нормировка производного и
потенциального членов не выведены из единого действия, наличие
майорановского конденсата не является предсказанием программы.
\end{nogocriterion}

\section{Запрет выбора ориентации чётным действием}

Две нижние ветви имеют ковариантные импульсы
\[
 k_{0}=+1,\qquad k_{-1}=-1
\]
в единичной геометрии. Для них совпадают все чётные спектральные
инварианты:
\[
 k_0^2=k_{-1}^2,\qquad
 k_0^4=k_{-1}^4,\qquad
 |k_0|=|k_{-1}|.
\]
Значит, всякое вещественное действие, зависящее только от \(D^2\) и чётных
степеней ковариантного импульса, сохраняет точное вырождение ветвей.

Ранее проверенная \(\eta\)-фаза в минимальной вектороподобной модели
сокращается между сопряжёнными блоками и не создаёт вещественного
расщепления энергии. Поэтому для выбора ориентации требуется либо
выведенный киральный граничный член, либо другой вещественный вклад,
нечётный относительно \(k\mapsto-k\).

\begin{proposition}[Двойная незамкнутость майорановской ветви]
Текущее спектральное ядро не определяет ни сам факт конденсации, ни выбор
одной из двух ориентаций. Первое зависит от невыведенной спектральной
функции и её масштаба, второе запрещено для вещественного чётного
спектрального действия.
\end{proposition}

\begin{protocol}[Условия продолжения]
\begin{enumerate}
 \item вывести спектральную функцию и её масштаб до обращения к
 нейтринным данным;
 \item получить коэффициенты производного, квадратичного и четвертичного
 членов из одного следа;
 \item доказать превышение геометрического порога;
 \item вывести вещественный ориентационно-нечётный член;
 \item после этого повторно вычислить ядро оператора Боголюбова--де Жена.
\end{enumerate}
\end{protocol}

Машинный протокол сохранён в
\texttt{s2t\_spectral\_pairing\_stiffness\_gate\_results.json}.